\documentclass[11pt]{article}
\usepackage[margin=1in]{geometry}
\usepackage{booktabs}
\usepackage{graphicx}
\usepackage{hyperref}
\usepackage{microtype}
\usepackage{enumitem}
\title{FinDS-AgentBench: A Pilot Benchmark for Financial Data-Science Agents}
\author{FinDS-AgentBench Contributors}
\date{}

\begin{document}
\maketitle

\begin{abstract}
We introduce FinDS-AgentBench, a benchmark for evaluating AI agents that perform end-to-end financial data-science research under temporal, leakage, reproducibility, and decision-quality constraints. The current pilot release contains 9 tasks, 9 runnable tasks, and 3 tracks. The repeated-run protocol covers 8 tasks and 35 task-system cells with fixed seeds and release-owned artifacts. In the pilot overall-score paired comparison against the best baseline per task, agents are higher in 0 cases, baselines are higher in 6 cases, and 2 cases are ties. The release ships task cards, data manifests, reference results, uncertainty tables, manual-audit artifacts, and a deterministic release reproducibility check.
\end{abstract}

\section{Introduction}
Financial machine-learning research is unusually sensitive to timestamp discipline, point-in-time data availability, target leakage, validation design, transaction costs, and narrative overclaiming. Generic data-science and tool-use benchmarks do not isolate these failure modes. FinDS-AgentBench is designed to evaluate whether agents can produce auditable financial research artifacts rather than merely plausible notebooks.

\paragraph{Contributions.}
\begin{enumerate}[leftmargin=*]
\item A pilot benchmark release with task specifications, public data manifests, task cards, and evaluation cards.
\item A command-based agent and baseline protocol with repeated runs, fixed seeds, clean release builds, and manifest validation.
\item Publication-facing reference results, uncertainty estimates, paired agent-vs-baseline comparisons, and paper-ready tables.
\item A manual-audit workflow covering temporal protocol correctness, leakage awareness, evidence use, claim discipline, and reproducibility trace completeness.
\end{enumerate}

\section{Related Work}
FinDS-AgentBench sits at the intersection of agent benchmarks, machine-learning experimentation benchmarks, finance-domain LLM evaluation, and leakage-aware temporal validation. General agent and software-engineering benchmarks such as ~\cite{agentbench2023,swebench2023} emphasize multi-step execution, but do not target financial ML research validity.
ML-oriented agent benchmarks such as ~\cite{mlagentbench2023,mlebench2024,scienceagentbench2024,mlrbench2025} evaluate experimentation, engineering, or open-ended research automation. FinDS-AgentBench adopts their emphasis on executable artifacts while adding point-in-time finance data, temporal holdouts, leakage scanning, and writeup audit.
Finance benchmarks such as ~\cite{financebench2023,finben2024,profitmirage2025,fintoolbench2026,finmcpbench2026,workstreambench2026,bluefin2026} cover financial QA, broad financial LLM capabilities, leakage in financial agents, tool use, and spreadsheet workflows. These are adjacent but do not provide an end-to-end financial ML research benchmark with point-in-time task cards, private temporal labels, repeated-run uncertainty, reproducibility gates, and narrative-discipline scoring.
Data-science and data-analysis benchmarks such as ~\cite{ds10002022,infiagentdabench2024} motivate reliable code execution and analysis-task evaluation. FinDS-AgentBench extends that execution focus to financial research protocols where model validation, data timestamp availability, and claims about economic usefulness must be evaluated jointly.
Temporal-reasoning benchmarks such as ~\cite{temporalbench2026} motivate time-aware evaluation. FinDS-AgentBench narrows that lens to finance, where point-in-time information availability and leakage-safe validation are central rather than incidental constraints.

\begin{table}[t]
\centering
\small
\begin{tabular}{llll}
\toprule
Work & Focus & Overlap & FinDS-AgentBench gap \\
\midrule
AgentBench~\cite{agentbench2023} & Interactive LLM-agent evaluation across multiple environments. & Frames LLMs as agents acting over multi-turn environments. & Not specialized for financial ML research artifacts, point-in-time data, leakage checks, or writeup validity. \\
SWE-bench~\cite{swebench2023} & Repository-level issue resolution from real GitHub issues and pull requests. & Execution-based, artifact-producing agent evaluation. & Software engineering rather than financial data-science research with temporal holdouts and financial validity gates. \\
DS-1000~\cite{ds10002022} & Data-science code generation with reliable functional tests. & Data analysis and modeling tasks with executable evaluation. & Short code-generation tasks rather than full financial research workflows with notebooks, writeups, and temporal validity checks. \\
InfiAgent-DABench~\cite{infiagentdabench2024} & Agentic data-analysis tasks over CSV datasets and execution environments. & End-to-end data analysis agents with executable code and structured evaluation. & General data analysis rather than finance-specific temporal information sets, private holdouts, leakage gates, and research writeup audit. \\
MLAgentBench~\cite{mlagentbench2023} & Language agents conducting ML experimentation tasks. & Long-horizon ML experimentation with file editing and code execution. & Not domain-specific to financial ML, point-in-time data availability, backtest realism, or claim audit. \\
MLE-bench~\cite{mlebench2024} & Kaggle-style ML engineering competitions for agents. & Agentic ML engineering, leaderboard-style scoring, and contamination analysis. & Broad ML engineering rather than finance-specific temporal information sets, leakage-safe research notebooks, and validity-gated writeups. \\
ScienceAgentBench~\cite{scienceagentbench2024} & Scientific workflow tasks extracted from peer-reviewed publications. & Research-style agent evaluation with code execution and expert validation. & Not specialized for market data, point-in-time financial constraints, financial metrics, or investment-style claim discipline. \\
MLR-Bench~\cite{mlrbench2025} & Open-ended ML research tasks with proposal, experimentation, and paper-writing stages. & Evaluates research automation and paper-style outputs. & Not focused on finance-specific leakage, temporal validation, private holdouts, or financial model-risk framing. \\
FinanceBench~\cite{financebench2023} & Open-book financial QA over company documents. & Finance-domain LLM evaluation with evidence requirements. & QA rather than executable financial ML research with notebooks, temporal holdouts, and private scoring. \\
FinBen~\cite{finben2024} & Broad financial LLM benchmark across language, QA, forecasting, risk, and decision tasks. & Finance-specific LLM evaluation across heterogeneous task types. & Broad LLM capability suite rather than validity-gated end-to-end financial ML research artifacts. \\
Profit Mirage / FinLake-Bench~\cite{profitmirage2025} & Information leakage in LLM-based financial agents and leakage-robust evaluation. & Closest finance-agent neighbor for leakage-aware evaluation. & Focuses on leakage and trading-agent generalization; FinDS-AgentBench broadens to reproducible ML research artifacts, task/evaluation cards, writeup audit, and multi-track validity scoring. \\
TemporalBench~\cite{temporalbench2026} & Contextual and event-informed time-series reasoning across multiple domains. & Temporal reasoning and event-conditioned prediction. & Multi-domain time-series reasoning rather than financial ML research artifacts with point-in-time market data, leakage scanning, and narrative audit. \\
FinToolBench~\cite{fintoolbench2026} & Financial tool-learning agents over executable financial tools. & Auditable financial agent execution and finance-specific tool use. & Tool invocation rather than leakage-safe financial ML research notebooks and validity-gated result claims. \\
FinMCP-Bench~\cite{finmcpbench2026} & Financial tool invocation through model context protocol servers. & Financial agent execution with real and synthetic tool-use queries. & MCP/tool-use benchmark rather than end-to-end financial ML research with temporal holdouts and artifact audit. \\
WorkstreamBench~\cite{workstreambench2026} & End-to-end financial spreadsheet construction and review. & Finance workflow artifacts and multidimensional professional-quality evaluation. & Spreadsheet modeling rather than financial ML research with temporal prediction tasks, private holdouts, and leakage/reproducibility gates. \\
BlueFin~\cite{bluefin2026} & Financial spreadsheet synthesis, manipulation, and comprehension tasks. & Finance-domain agent artifacts with granular rubric-based evaluation. & Spreadsheet workbook workflows rather than executable financial ML research notebooks, temporal prediction tasks, and leakage-safe private scoring. \\
\bottomrule
\end{tabular}
\caption{Related benchmark positioning. FinDS-AgentBench targets the missing intersection of end-to-end financial ML research artifacts, temporal information discipline, leakage resistance, reproducibility, and writeup validity.}
\label{tab:related-work-positioning}
\end{table}


\section{Benchmark Design}
The pilot release spans event\_aware\_time\_series\_reasoning, predictive\_financial\_ml, research\_replication\_and\_audit. Each task specification identifies the prediction timestamp, allowed public information, target horizon, temporal split, required submission artifacts, primary metric, and leakage constraints. The benchmark separates predictive performance from artifact validity so that a high private-holdout score does not erase reproducibility, leakage, or methodology failures.

\section{Pilot Release}
The release manifest records 9 tasks and 9 runnable tasks. The repeated-run pilot protocol covers 8 tasks, excluding the leakage-audit task from the predictive baseline/agent sweep because it uses an audit-style scoring surface rather than a standard prediction holdout.

\section{Evaluation Protocol}
The combined pilot protocol evaluates 27 baseline task-system cells and 8 agent task-system cells. Observed run counts per cell are 3. Baselines include naive and classical models; the bundled example agents are environment-wrapped systems that select among public candidate strategies before writing holdout predictions and research artifacts. The external-agent readiness gate is \texttt{not\_ready\_no\_external\_agents} with 0 completed external agent configurations.

\paragraph{Statistical reporting.}
We report repeated-run uncertainty from the combined protocol run table \texttt{reports/release\_runs/pilot\_protocol\_run\_results.csv}. The reported metrics are score.overall\_score, score.balanced\_accuracy, score.roc\_auc.
For each task, system, run type, and metric, completed runs with numeric metric values are summarized by the sample mean, sample standard deviation, standard error, and a two-sided 95\% t-interval. Single-run groups are reported with zero-width intervals rather than extrapolated uncertainty.
For agent-vs-baseline comparisons, the best baseline is selected independently for each task and metric using the completed-run mean. Agent and baseline runs are paired by \texttt{trace.repeat\_index}, falling back to \texttt{trace.seed} only when no repeat index is available. We report the paired agent-minus-baseline mean difference, its t-interval, and an exact two-sided sign-test p-value after dropping zero differences.
Because the pilot uses a small repeated-run count, these statistics are treated as transparent uncertainty and regression artifacts rather than definitive significance claims.


\section{Results}
Table~\ref{tab:pilot-agent-vs-best-baseline-score-overall-score} summarizes the paired overall-score comparison between each agent and the best completed-run baseline for the same task. Appendix Table~\ref{tab:pilot-uncertainty-score-overall-score} reports overall-score repeated-run uncertainty for all task-system cells.

\begin{table}[t]
\centering
\small
\begin{tabular}{lllrrrrr}
\toprule
Task & Agent & Best Baseline & $n$ & Agent & Baseline & $\Delta$ & $p_{sign}$ \\
\midrule
Front-End Spread & Front-End Sweep Env-Agent & LogReg & 3 & 0.543 & 0.543 & 0.000 &  \\
Synthetic Event & Event Rule Env-Agent & Event Rule & 3 & 0.637 & 0.637 & 0.000 &  \\
Synthetic Market & Market Sweep Env-Agent & Momentum & 3 & 0.511 & 0.521 & -0.010 & 0.250 \\
USD AFE vs. EME & USD Sweep Env-Agent & Extra Trees & 3 & 0.507 & 0.529 & -0.022 & 0.250 \\
USD Broad & USD Sweep Env-Agent & Prev-Day & 3 & 0.502 & 0.510 & -0.008 & 0.250 \\
Treasury 10Y-2Y Curve & Treasury 10Y-2Y Sweep Env-Agent & LogReg & 3 & 0.501 & 0.530 & -0.029 & 0.250 \\
Treasury 10Y-3M Curve & Treasury 10Y-3M Sweep Env-Agent & Extra Trees & 3 & 0.517 & 0.533 & -0.016 & 0.250 \\
Treasury 10Y Direction & Treasury 10Y Sweep Env-Agent & Extra Trees & 3 & 0.504 & 0.506 & -0.002 & 1.000 \\
\bottomrule
\end{tabular}
\caption{Paired agent-minus-best-baseline comparison for Overall. Baselines are selected separately within each task by completed-run mean. The sign-test p-value is exact and two-sided.}
\label{tab:pilot-agent-vs-best-baseline-score-overall-score}
\end{table}


\section{Manual Audit and Validity Checks}
The pilot manual-audit bundle currently includes 6 cases across 3 reviewed tasks. Its official status is \texttt{seed\_author\_adjudication\_only}; independent-overlap agreement status is \texttt{insufficient\_independent\_overlap}, with exploratory status \texttt{pairwise\_agreement\_available}. The reviewer-readiness gate is \texttt{not\_ready\_seed\_only} with 0 completed independent reviewer packets. This makes the current audit layer suitable as a seed adjudication workflow, while a submission-ready paper should add an independent second-reviewer packet.

\section{Qualitative Failure Examples}
The manual audit illustrates why predictive scores and automatic artifact validators are insufficient for financial research-agent evaluation. The examples below are generated from the seed adjudicated audit subset and preserve task, run, and artifact references.

\begin{description}[leftmargin=*]
\item[pilot\_event\_rule\_baseline\_release\_001] \textbf{Task/system:} \texttt{synthetic\_event\_response\_v0} / \texttt{event\_rule\_baseline} (baseline, run label \texttt{release\_pilot\_event\_001\_seed\_23}). \textbf{Audit label:} \texttt{minimally\_defensible}, total score 6.
\textbf{Low-scoring dimensions:} temporal\_protocol\_correctness, quantitative\_evidence\_use, baseline\_comparison\_or\_counterfactual\_context, reproducibility\_trace\_completeness. \textbf{Primary findings:} Strong predictive score is not accompanied by quantitative narrative support; Temporal protocol description is brief; No comparator or ablation context appears in the writeup. \textbf{Evidence:} No ablation, heuristic comparison, or simpler event rule is discussed.
\textbf{Artifact reference:} \texttt{runs/suites/pilot\_baselines\_v0/synthetic\_event\_response\_v0/event\_rule\_baseline/release\_pilot\_event\_001\_seed\_23/writeup.md}.

\item[pilot\_market\_momentum\_baseline\_release\_001] \textbf{Task/system:} \texttt{synthetic\_market\_direction\_v0} / \texttt{momentum\_baseline} (baseline, run label \texttt{release\_pilot\_market\_001\_seed\_11}). \textbf{Audit label:} \texttt{minimally\_defensible}, total score 6.
\textbf{Low-scoring dimensions:} temporal\_protocol\_correctness, quantitative\_evidence\_use, baseline\_comparison\_or\_counterfactual\_context, reproducibility\_trace\_completeness. \textbf{Primary findings:} No quantitative support in the writeup; No baseline or counterfactual comparison; Temporal protocol rationale is thinner than the artifact package itself. \textbf{Evidence:} The submission does not compare the momentum rule to any simpler or more expressive alternative.
\textbf{Artifact reference:} \texttt{runs/suites/pilot\_baselines\_v0/synthetic\_market\_direction\_v0/momentum\_baseline/release\_pilot\_market\_001\_seed\_11/writeup.md}.

\item[pilot\_event\_rule\_env\_agent\_protocol\_001] \textbf{Task/system:} \texttt{synthetic\_event\_response\_v0} / \texttt{event\_rule\_env\_agent} (agent, run label \texttt{release\_protocol\_agent\_event\_001\_seed\_23}). \textbf{Audit label:} \texttt{minimally\_defensible}, total score 7.
\textbf{Low-scoring dimensions:} temporal\_protocol\_correctness, quantitative\_evidence\_use, baseline\_comparison\_or\_counterfactual\_context. \textbf{Primary findings:} Agent trace provenance is stronger than the baseline trace; Narrative quality is still thin; No quantitative or comparative justification appears in the writeup. \textbf{Evidence:} The agent output does not situate the rule against any competing event heuristic.
\textbf{Artifact reference:} \texttt{runs/suites/pilot\_protocol\_v0/synthetic\_event\_response\_v0/event\_rule\_env\_agent/release\_protocol\_agent\_event\_001\_seed\_23/writeup.md}.

\item[pilot\_market\_logistic\_baseline\_release\_001] \textbf{Task/system:} \texttt{synthetic\_market\_direction\_v0} / \texttt{logistic\_regression\_baseline} (baseline, run label \texttt{release\_pilot\_market\_001\_seed\_11}). \textbf{Audit label:} \texttt{strong\_but\_incomplete}, total score 10.
\textbf{Low-scoring dimensions:} baseline\_comparison\_or\_counterfactual\_context. \textbf{Primary findings:} Method description is strong; Quantitative support is present; Main remaining gap is missing comparison against simpler alternatives. \textbf{Evidence:} Even though a momentum baseline exists in the suite, the writeup does not compare against it or justify model complexity relative to a rule baseline.
\textbf{Artifact reference:} \texttt{runs/suites/pilot\_baselines\_v0/synthetic\_market\_direction\_v0/logistic\_regression\_baseline/release\_pilot\_market\_001\_seed\_11/writeup.md}.

\end{description}


\section{Reproducibility and Release Artifacts}
The release build is intentionally scriptable. It rebuilds suite reports, reference results, statistical artifacts, paper artifacts, task cards, data manifests, and the canonical release manifest. A separate smoke check builds the pilot release twice in isolated roots and compares deterministic publication-facing outputs by content digest.

\begin{verbatim}
PYTHONPATH=src python scripts/build_pilot_release.py --repeat 3 --market-seed 11 --event-seed 23 --treasury-seed 29 --curve-seed 31 --curve3mo-seed 33 --front-end-seed 31 --usd-seed 37 --treasury-snapshot-date 2026-06-21 --curve-snapshot-date 2026-06-21 --curve3mo-snapshot-date 2026-06-21 --front-end-snapshot-date 2026-06-21 --usd-snapshot-date 2026-06-21 --clean-existing-outputs
\end{verbatim}

\section{Limitations}
\begin{itemize}[leftmargin=*]
\item The current pilot uses a small repeated-run count, so uncertainty artifacts should be read as transparency and regression checks rather than definitive significance claims.
\item Bundled agents are controlled example wrappers, not a broad external-agent leaderboard.
\item The manual-audit layer still needs independent second-reviewer completion before submission-strength claims.
\item The task set is strong enough for a workshop pilot, but a top benchmark/dataset venue version should scale to a larger frozen suite with hidden temporal holdouts.
\end{itemize}

\section{Conclusion}
FinDS-AgentBench frames financial data-science agent evaluation as a joint test of predictive performance, temporal validity, leakage resistance, reproducibility, and claim discipline. The pilot release provides the first reproducible scaffold for this evaluation and establishes the artifacts needed for an arXiv/workshop submission.

\appendix
\section{Full Repeated-Run Uncertainty Table}
\begin{table}[t]
\centering
\small
\begin{tabular}{llllcc}
\toprule
Task & System & Type & $n$ & Mean & 95\% CI \\
\midrule
Front-End Spread & Extra Trees & Baseline & 3 & 0.509 & [0.475, 0.542] \\
Front-End Spread & LogReg & Baseline & 3 & 0.543 & [0.543, 0.543] \\
Front-End Spread & Prev-Day & Baseline & 3 & 0.540 & [0.540, 0.540] \\
Front-End Spread & Rand. Forest & Baseline & 3 & 0.511 & [0.492, 0.529] \\
Front-End Spread & Front-End Sweep Env-Agent & Agent & 3 & 0.543 & [0.543, 0.543] \\
Synthetic Event & Event Rule & Baseline & 3 & 0.637 & [0.611, 0.664] \\
Synthetic Event & Event Rule Env-Agent & Agent & 3 & 0.637 & [0.611, 0.664] \\
Synthetic Market & LogReg & Baseline & 3 & 0.514 & [0.475, 0.554] \\
Synthetic Market & Market Sweep Env-Agent & Agent & 3 & 0.511 & [0.461, 0.562] \\
Synthetic Market & Momentum & Baseline & 3 & 0.521 & [0.475, 0.567] \\
USD AFE vs. EME & Extra Trees & Baseline & 3 & 0.529 & [0.523, 0.535] \\
USD AFE vs. EME & LogReg & Baseline & 3 & 0.500 & [0.500, 0.500] \\
USD AFE vs. EME & Prev-Day & Baseline & 3 & 0.511 & [0.511, 0.511] \\
USD AFE vs. EME & Rand. Forest & Baseline & 3 & 0.512 & [0.481, 0.543] \\
USD AFE vs. EME & USD Sweep Env-Agent & Agent & 3 & 0.507 & [0.490, 0.524] \\
USD Broad & Extra Trees & Baseline & 3 & 0.500 & [0.493, 0.506] \\
USD Broad & LogReg & Baseline & 3 & 0.507 & [0.507, 0.507] \\
USD Broad & Prev-Day & Baseline & 3 & 0.510 & [0.510, 0.510] \\
USD Broad & Rand. Forest & Baseline & 3 & 0.500 & [0.485, 0.515] \\
USD Broad & USD Sweep Env-Agent & Agent & 3 & 0.502 & [0.488, 0.515] \\
Treasury 10Y-2Y Curve & Extra Trees & Baseline & 3 & 0.501 & [0.492, 0.510] \\
Treasury 10Y-2Y Curve & LogReg & Baseline & 3 & 0.530 & [0.530, 0.530] \\
Treasury 10Y-2Y Curve & Prev-Day & Baseline & 3 & 0.500 & [0.500, 0.500] \\
Treasury 10Y-2Y Curve & Rand. Forest & Baseline & 3 & 0.498 & [0.489, 0.506] \\
Treasury 10Y-2Y Curve & Treasury 10Y-2Y Sweep Env-Agent & Agent & 3 & 0.501 & [0.492, 0.510] \\
Treasury 10Y-3M Curve & Extra Trees & Baseline & 3 & 0.533 & [0.505, 0.561] \\
Treasury 10Y-3M Curve & LogReg & Baseline & 3 & 0.517 & [0.517, 0.517] \\
Treasury 10Y-3M Curve & Prev-Day & Baseline & 3 & 0.500 & [0.500, 0.500] \\
Treasury 10Y-3M Curve & Rand. Forest & Baseline & 3 & 0.532 & [0.513, 0.552] \\
Treasury 10Y-3M Curve & Treasury 10Y-3M Sweep Env-Agent & Agent & 3 & 0.517 & [0.517, 0.517] \\
Treasury 10Y Direction & Extra Trees & Baseline & 3 & 0.506 & [0.487, 0.526] \\
Treasury 10Y Direction & LogReg & Baseline & 3 & 0.487 & [0.487, 0.487] \\
Treasury 10Y Direction & Prev-Day & Baseline & 3 & 0.494 & [0.494, 0.494] \\
Treasury 10Y Direction & Rand. Forest & Baseline & 3 & 0.503 & [0.482, 0.525] \\
Treasury 10Y Direction & Treasury 10Y Sweep Env-Agent & Agent & 3 & 0.504 & [0.494, 0.514] \\
\bottomrule
\end{tabular}
\caption{Pilot repeated-run uncertainty for Overall. Intervals are t-based 95\% confidence intervals over completed runs.}
\label{tab:pilot-uncertainty-score-overall-score}
\end{table}


\section{Full Pilot Protocol Reference Table}
\begin{table}[t]
\centering
\small
\resizebox{\textwidth}{!}{%
\begin{tabular}{lllcccc}
\toprule
Task & Agent & Type & $n$ & Overall & Balanced Acc. & ROC AUC \\
\midrule
Front-End Spread & Extra Trees & Baseline & 3 & 0.509 \pm 0.013 & 0.509 \pm 0.013 & 0.529 \pm 0.003 \\
Front-End Spread & LogReg & Baseline & 3 & 0.543 \pm 0.000 & 0.543 \pm 0.000 & 0.536 \pm 0.000 \\
Front-End Spread & Prev-Day & Baseline & 3 & 0.540 \pm 0.000 & 0.540 \pm 0.000 & 0.540 \pm 0.000 \\
Front-End Spread & Rand. Forest & Baseline & 3 & 0.511 \pm 0.008 & 0.511 \pm 0.008 & 0.518 \pm 0.003 \\
Front-End Spread & Front-End Sweep Env-Agent & Agent & 3 & 0.543 \pm 0.000 & 0.543 \pm 0.000 & 0.536 \pm 0.000 \\
Synthetic Event & Event Rule & Baseline & 3 & 0.637 \pm 0.011 & 0.637 \pm 0.011 & 0.691 \pm 0.018 \\
Synthetic Event & Event Rule Env-Agent & Agent & 3 & 0.637 \pm 0.011 & 0.637 \pm 0.011 & 0.691 \pm 0.018 \\
Synthetic Market & LogReg & Baseline & 3 & 0.514 \pm 0.016 & 0.514 \pm 0.016 & 0.528 \pm 0.028 \\
Synthetic Market & Momentum & Baseline & 3 & 0.521 \pm 0.018 & 0.521 \pm 0.018 & 0.518 \pm 0.024 \\
Synthetic Market & Market Sweep Env-Agent & Agent & 3 & 0.511 \pm 0.020 & 0.511 \pm 0.020 & 0.518 \pm 0.023 \\
USD AFE vs. EME & Extra Trees & Baseline & 3 & 0.529 \pm 0.002 & 0.529 \pm 0.002 & 0.541 \pm 0.002 \\
USD AFE vs. EME & LogReg & Baseline & 3 & 0.500 \pm 0.000 & 0.500 \pm 0.000 & 0.542 \pm 0.000 \\
USD AFE vs. EME & Prev-Day & Baseline & 3 & 0.511 \pm 0.000 & 0.511 \pm 0.000 & 0.511 \pm 0.000 \\
USD AFE vs. EME & Rand. Forest & Baseline & 3 & 0.512 \pm 0.013 & 0.512 \pm 0.013 & 0.532 \pm 0.002 \\
USD AFE vs. EME & USD Sweep Env-Agent & Agent & 3 & 0.507 \pm 0.007 & 0.507 \pm 0.007 & 0.517 \pm 0.011 \\
USD Broad & Extra Trees & Baseline & 3 & 0.500 \pm 0.003 & 0.500 \pm 0.003 & 0.490 \pm 0.007 \\
USD Broad & LogReg & Baseline & 3 & 0.507 \pm 0.000 & 0.507 \pm 0.000 & 0.495 \pm 0.000 \\
USD Broad & Prev-Day & Baseline & 3 & 0.510 \pm 0.000 & 0.510 \pm 0.000 & 0.510 \pm 0.000 \\
USD Broad & Rand. Forest & Baseline & 3 & 0.500 \pm 0.006 & 0.500 \pm 0.006 & 0.506 \pm 0.006 \\
USD Broad & USD Sweep Env-Agent & Agent & 3 & 0.502 \pm 0.005 & 0.502 \pm 0.005 & 0.497 \pm 0.015 \\
Treasury 10Y-2Y Curve & Extra Trees & Baseline & 3 & 0.501 \pm 0.004 & 0.501 \pm 0.004 & 0.522 \pm 0.000 \\
Treasury 10Y-2Y Curve & LogReg & Baseline & 3 & 0.530 \pm 0.000 & 0.530 \pm 0.000 & 0.558 \pm 0.000 \\
Treasury 10Y-2Y Curve & Prev-Day & Baseline & 3 & 0.500 \pm 0.000 & 0.500 \pm 0.000 & 0.530 \pm 0.000 \\
Treasury 10Y-2Y Curve & Rand. Forest & Baseline & 3 & 0.498 \pm 0.004 & 0.498 \pm 0.004 & 0.511 \pm 0.001 \\
Treasury 10Y-2Y Curve & Treasury 10Y-2Y Sweep Env-Agent & Agent & 3 & 0.501 \pm 0.004 & 0.501 \pm 0.004 & 0.522 \pm 0.000 \\
Treasury 10Y-3M Curve & Extra Trees & Baseline & 3 & 0.533 \pm 0.011 & 0.533 \pm 0.011 & 0.550 \pm 0.003 \\
Treasury 10Y-3M Curve & LogReg & Baseline & 3 & 0.517 \pm 0.000 & 0.517 \pm 0.000 & 0.526 \pm 0.000 \\
Treasury 10Y-3M Curve & Prev-Day & Baseline & 3 & 0.500 \pm 0.000 & 0.500 \pm 0.000 & 0.511 \pm 0.000 \\
Treasury 10Y-3M Curve & Rand. Forest & Baseline & 3 & 0.532 \pm 0.008 & 0.532 \pm 0.008 & 0.540 \pm 0.006 \\
Treasury 10Y-3M Curve & Treasury 10Y-3M Sweep Env-Agent & Agent & 3 & 0.517 \pm 0.000 & 0.517 \pm 0.000 & 0.526 \pm 0.000 \\
Treasury 10Y Direction & Extra Trees & Baseline & 3 & 0.506 \pm 0.008 & 0.506 \pm 0.008 & 0.525 \pm 0.003 \\
Treasury 10Y Direction & LogReg & Baseline & 3 & 0.487 \pm 0.000 & 0.487 \pm 0.000 & 0.486 \pm 0.000 \\
Treasury 10Y Direction & Prev-Day & Baseline & 3 & 0.494 \pm 0.000 & 0.494 \pm 0.000 & 0.494 \pm 0.000 \\
Treasury 10Y Direction & Rand. Forest & Baseline & 3 & 0.503 \pm 0.009 & 0.503 \pm 0.009 & 0.523 \pm 0.004 \\
Treasury 10Y Direction & Treasury 10Y Sweep Env-Agent & Agent & 3 & 0.504 \pm 0.004 & 0.504 \pm 0.004 & 0.527 \pm 0.000 \\
\bottomrule
\end{tabular}
}
\caption{Combined Pilot Protocol repeated-run summary. Means and sample standard deviations are computed over the configured seeds.}
\label{tab:pilot-protocol}
\end{table}


\bibliographystyle{plain}
\bibliography{references}

\end{document}
